%!TEX root = tfg.tex
\chapter{Introducción}

\begin{abstract}
Este capítulo presenta una visión general de \textit{Via Inferni} y sitúa el proyecto
dentro del marco de un Trabajo de Fin de Grado de Ingeniería del Software. Se resume
la propuesta del videojuego, el problema que aborda y la forma en la que se estructura
la memoria tras la reorganización acordada.
\end{abstract}

\section{Presentación del proyecto}
\textit{Via Inferni} es un videojuego \textit{roguelike} 2D con generación procedural
de niveles inspirado en la estructura del Infierno de Dante. El proyecto combina una
base técnica propia del desarrollo de videojuegos independientes con una propuesta
narrativa clásica, articulando el progreso del jugador a través de distintos círculos,
enemigos temáticos y sistemas de combate, exploración y progresión.

Desde el punto de vista académico, este TFG se plantea como un ejercicio integral de
ingeniería del software. No solo se aborda la implementación del producto, sino también
su planificación, la gestión del alcance, la organización del trabajo por funcionalidades,
el control del código y la documentación del proceso seguido hasta obtener una versión
jugable del sistema.

\section{Estructura de la memoria}
La memoria se ha organizado en cinco grandes bloques. El primero, \textbf{Inicio},
recoge el contexto del proyecto, sus objetivos y la metodología utilizada. El segundo,
\textbf{Planificación}, desarrolla los planes de gestión y explica cómo se ejecutaron
durante el proyecto. El tercer bloque, \textbf{Ejecución, seguimiento y control},
documenta el desarrollo iterativo del videojuego. El cuarto, \textbf{Cierre}, agrupa
los materiales finales de uso, instalación y conclusiones. Por último, los
\textbf{Anexos} recogen los glosarios y las referencias bibliográficas.
